\section{Controle com reação nula e manipuladores \textit{free-floating}}
\label{sec_reaction_null_free_floating}

Nesta seção é apresentada uma revisão da literatura relacionada ao controle de manipuladores espaciais \textit{free-floating} com redução ou anulação das reações transmitidas à base. São destacadas as referências que tratam do conceito de \textit{reaction null space}, do controle \textit{reactionless}, da captura de alvos com mínima perturbação da base e de estratégias adaptativas aplicadas a sistemas espaciais acoplados.

Entre as contribuições específicas da literatura de robótica espacial, destaca-se a linha de pesquisa voltada ao controle de manipuladores \textit{free-floating} com redução ou anulação das reações dinâmicas impostas à base. Nessa classe de problemas, o objetivo não é apenas garantir o rastreamento da trajetória do efetuador terminal, mas fazê-lo de modo a preservar a atitude da espaçonave ou a integridade dinâmica da estrutura que suporta o manipulador. Essa questão é particularmente importante em missões de proximidade, captura e serviço em órbita, nas quais perturbações excessivas da base podem comprometer a segurança da operação e aumentar o consumo de atuadores auxiliares.

Um dos trabalhos mais importantes nessa linha é apresentado em \cite{nenchev1999reaction}, onde é proposta uma lei de controle composta para o rastreamento da trajetória do efetuador terminal em um sistema manipulador montado sobre estrutura flexível, de forma a não induzir perturbações na base. Os autores formulam o problema com base no conceito de \textit{reaction null space}, introduzido para tratar a interação dinâmica em robôs \textit{free-floating} e sistemas de base móvel em geral.
Além da reação nula, o trabalho inclui a supressão de vibração da base e apresenta validação por simulação e por meio do banco experimental, desenvolvido na \gls{trep}.
Essa referência tornou-se importante por consolidar o conceito de \textit{reaction null space} como ferramenta de projeto para desacoplar, ao menos parcialmente, o movimento útil do manipulador das perturbações transmitidas à estrutura de suporte.

\cite{piersigilli2010reactionless} investigam o problema da captura \textit{reactionless} de um satélite por um manipulador de dois graus de liberdade. Nesse trabalho, os autores tratam explicitamente a captura sem reação de um alvo por um manipulador espacial, reforçando a aplicação prática desse tipo de estratégia em operações de contato e apreensão. Essa referência é relevante porque desloca a discussão de reação nula do campo conceitual para uma tarefa operacional crítica, na qual impulsos e perturbações na base tornam-se ainda mais sensíveis.

Para ampliar essa discussão, \cite{nguyenhuynh2013adaptive} consideram o problema de movimento \textit{reactionless} adaptativo e identificação de parâmetros após a captura de detritos espaciais. Os autores mostram que a captura de um alvo altera a dinâmica do sistema e introduz incertezas paramétricas adicionais, tornando insuficientes estratégias baseadas em modelos exatamente conhecidos. Dessa forma, a combinação entre redução de reação e identificação paramétrica passa a ocupar papel importante na literatura, aproximando o problema de cenários mais realistas de serviço em órbita e pós-captura.

\cite{xu2016adaptive} apresentam um controle adaptativo de movimento com reação nula para manipuladores espaciais \textit{free-floating} com incertezas cinemáticas e dinâmicas. Os autores destacam que uma das principais dificuldades para a construção de uma lei adaptativa baseada em \textit{reaction null space} está na obtenção de uma expressão linear em relação aos parâmetros incertos. A partir dessa reformulação, é proposto um controlador em nível de velocidade capaz de regular a atitude da espaçonave e realizar o rastreamento da trajetória do efetuador terminal. O método é demonstrado por simulações numéricas em um manipulador planar de três graus de liberdade. Essa referência mostra um avanço da literatura ao levar a ideia de reação nula para um cenário com incertezas mais próximas das encontradas em missões reais.

Mais recentemente, \cite{ye2019adaptive} apresentam uma estratégia adaptativa de planejamento e controle baseada em \textit{Adaptive Reaction Null Space} para robôs espaciais \textit{free-floating} em contato com alvos desconhecidos. O método combina um procedimento de planejamento adaptativo via \gls{vffrls} com uma estratégia robusta de controle baseada em \gls{ssadeelm}, buscando compensar distúrbios desconhecidos e incertezas do modelo sem exigir conhecimento preciso dos atributos do alvo.
Os autores apresentam comparações por simulação e experimento com métodos existentes, indicando efetividade e superioridade da estratégia proposta no cenário estudado. Essa referência mostra a evolução da literatura em direção a abordagens adaptativas voltadas a incertezas operacionais e interação com alvos reais.

De modo geral, a literatura revisada mostra que o problema do controle de manipuladores espaciais não se resume ao rastreamento geométrico de trajetórias, envolvendo também a gestão das reações dinâmicas impostas à base. Os trabalhos analisados mostram uma evolução que vai das formulações conceituais de \textit{reaction null space} e experimentos em estruturas flexíveis até estratégias adaptativas voltadas à captura, ao pós-captura e ao contato com alvos desconhecidos.
Esses aspectos são relevantes para a presente dissertação, pois reforçam que, em sistemas acoplados espaçonave--manipulador, o desempenho do braço robótico e o comportamento dinâmico da espaçonave não podem ser tratados de forma dissociada, especialmente nas fases de aproximação final e interação física com o alvo.